\chapter{Green's Functions}

\begin{remarkbox}
Green's functions are the response functions of linear differential equations.
They convert differential equations with sources into integral formulas. In field
theory they appear as retarded responses, advanced responses, static potentials,
and, after quantization, propagators.
\end{remarkbox}

\section{Green's Functions for Ordinary Differential Equations}

Consider a linear differential operator \(L\) acting on functions of one
variable. A Green's function \(G(x,y)\) satisfies
\[
    L_xG(x,y)=\delta(x-y),
\]
where \(L_x\) means that \(L\) acts on the \(x\)-dependence. If one wants to
solve
\[
    Lf(x)=s(x),
\]
then formally
\[
    f(x)=\int dy\,G(x,y)s(y)
\]
provided the boundary conditions are correctly encoded in \(G\).

The boundary conditions are part of the definition. The same differential
operator can have different Green's functions for Dirichlet, Neumann, periodic,
retarded, or advanced boundary conditions.

\begin{definitionbox}
A Green's function is an inverse kernel for a linear differential operator,
subject to specified boundary or causality conditions.
\end{definitionbox}

As a simple example, for
\[
    -\frac{d^2}{dx^2}G(x,y)=\delta(x-y)
\]
on an interval, the Green's function depends on what is required at the endpoints.
Thus it is misleading to speak of ``the'' Green's function without specifying the
problem.

\section{Green's Functions for Partial Differential Equations}

For a linear partial differential equation
\[
    L\phi(x)=J(x),
\]
a Green's function satisfies
\[
    L_xG(x,y)=\delta^{(D)}(x-y),
\]
where \(D=d+1\) is the spacetime dimension or \(D=d\) for a purely spatial
problem. The solution is formally
\[
    \phi(x)=\int d^D y\,G(x,y)J(y).
\]

If the operator is translation invariant, then
\[
    G(x,y)=G(x-y).
\]
In that case Fourier transform is often the most efficient tool. If
\[
    L(\partial)G(x)=\delta^{(D)}(x),
\]
then in momentum space one has
\[
    L(ik)\widetilde{G}(k)=1,
\]
so
\[
    \widetilde{G}(k)=\frac{1}{L(ik)}.
\]
Poles of \(\widetilde{G}(k)\) encode the characteristic modes of the differential
operator.

\begin{remarkbox}
The formal inverse \(1/L(ik)\) is not yet a full Green's function. One must also
specify how poles are treated. That prescription encodes causality or boundary
conditions.
\end{remarkbox}

\section{Retarded, Advanced, and Feynman-Type Green's Functions}

For hyperbolic equations, different Green's functions encode different causal
conditions. The retarded Green's function satisfies
\[
    G_{\mathrm{ret}}(x-y)=0
\]
unless \(x\) lies in or on the future light cone of \(y\). It gives causal
response to past sources:
\[
    \phi(x)=\int d^D y\,G_{\mathrm{ret}}(x-y)J(y).
\]

The advanced Green's function satisfies the opposite support condition. It is
nonzero only when \(x\) lies in or on the past light cone of \(y\). It is useful
mathematically but usually not used for causal initial-value evolution.

In quantum field theory, the Feynman propagator appears. It is not simply the
classical retarded Green's function. It implements time ordering and a specific
pole prescription. Formally, for a scalar field it is associated with
\[
    \frac{1}{k^2-m^2+i\epsilon}
\]
up to metric and Fourier conventions.

\begin{definitionbox}
Retarded Green's functions encode causal response. Advanced Green's functions
encode anti-causal response. Feynman-type Green's functions encode time-ordered
quantum propagation.
\end{definitionbox}

\section{Green's Functions for the Wave Equation}

The massless wave equation with source is
\[
    \Box\phi=J,
\]
where
\[
    \Box=\partial_t^2-\nabla^2.
\]
A Green's function satisfies
\[
    \Box G(x-y)=\delta^{(d+1)}(x-y).
\]
The retarded Green's function gives the field produced by sources in the past.
In three spatial dimensions, the retarded solution has support on the light cone.
This reflects sharp propagation at the speed of light.

For a source \(J\), the retarded solution is schematically
\[
    \phi_{\mathrm{ret}}(x)=\int d^{d+1}y\,G_{\mathrm{ret}}(x-y)J(y).
\]
If initial data are also present, the full solution is the sum of the homogeneous
solution determined by initial data and a particular solution determined by the
source.

\begin{remarkbox}
A Green's function solves the inhomogeneous problem with a point source. General
sources are handled by superposition, but only because the differential equation
is linear.
\end{remarkbox}

\section{Green's Functions for the Klein--Gordon Equation}

The sourced Klein--Gordon equation is
\[
    (\Box+m^2)\phi=J.
\]
A Green's function satisfies
\[
    (\Box_x+m^2)G(x-y)=\delta^{(d+1)}(x-y).
\]
In Fourier representation, one formally writes
\[
    G(x)=\int\frac{d^{d+1}k}{(2\pi)^{d+1}}\,
    \frac{e^{-ik_\mu x^\mu}}{-k_\mu k^\mu+m^2}
\]
with a pole prescription.

The massive case differs from the massless wave equation in two important ways.
First, the dispersion relation is
\[
    \omega^2=|\vect{k}|^2+m^2.
\]
Second, the static Green's function in three spatial dimensions decays
exponentially:
\[
    G_{\mathrm{static}}(r)\propto\frac{e^{-mr}}{r}.
\]
This is the Yukawa behavior. The length scale is
\[
    \ell_m=\frac{1}{m}.
\]

In the massless limit,
\[
    \frac{e^{-mr}}{r}\longrightarrow\frac{1}{r}.
\]
Thus a massive field has short-range static response, while a massless field has
long-range static response.

\section{Physical Interpretation as Response to Sources}

The most useful physical interpretation of a Green's function is response to a
point source. If
\[
    L_xG(x,y)=\delta(x-y),
\]
then \(G(x,y)\) is the field at \(x\) produced by a unit source placed at \(y\),
with the specified boundary or causal prescription.

For a general source \(J(y)\), linearity gives
\[
    \phi(x)=\int dy\,G(x,y)J(y).
\]
This says that the field is a superposition of responses to infinitesimal point
sources.

In electromagnetism, retarded Green's functions give retarded potentials. In
scalar field theory, they give causal response to scalar sources. In boundary
value problems, Green's functions encode how boundaries influence fields. In
quantum field theory, propagators encode how disturbances and correlations are
transmitted.

\begin{summarybox}
A Green's function is an inverse kernel for a linear differential operator plus a
choice of boundary or causal prescription. Retarded Green's functions describe
causal response, static Green's functions describe potentials, and the
Klein--Gordon Green's function displays the difference between massive
short-range response and massless long-range response.
\end{summarybox}
